\documentclass[11pt,a4paper]{article}
\usepackage[margin=1in]{geometry}
\usepackage{amsmath,amssymb,amsfonts,amsthm}
\usepackage{booktabs}
\usepackage{hyperref}

\newtheorem{proposition}{Proposition}
\newtheorem{lemma}{Lemma}

\title{Curvature Invariants of $\mathbb{C}P^2 \times S^3$\\
for the $a_4$ Seeley--DeWitt Coefficient\\[6pt]
\large DFD $\alpha$ Derivation: Supporting Calculations}

\author{Research Computation --- March 2026}
\date{}

\begin{document}
\maketitle

%=============================================================================
\section{Conventions and Setup}
%=============================================================================

We compute all curvature invariants needed for the $a_4$ Seeley--DeWitt
heat kernel coefficient on the internal space $X = \mathbb{C}P^2 \times S^3$
(real dimension~7) appearing in the DFD derivation of the fine structure
constant.

\medskip\noindent\textbf{Conventions:}
\begin{itemize}
\item $\mathbb{C}P^2$: Fubini--Study metric with holomorphic sectional
  curvature $c = 4$, so sectional curvatures lie in $[1,4]$.
  Real dimension $n = 4$.
\item $S^3$: Round metric of radius $r$. Constant sectional curvature
  $K = 1/r^2$. Real dimension $n = 3$.
\item Sign convention: $R^i{}_{jkl} = \partial_k \Gamma^i_{jl} - \cdots$
  with $\mathrm{Ric}_{jl} = R^i{}_{jil}$ (positive for spheres).
\item The $a_4$ coefficient follows Vassilevich~\cite{Vassilevich2003},
  Eq.~(4.3).
\end{itemize}

%=============================================================================
\section{$\mathbb{C}P^2$ with Fubini--Study Metric}
%=============================================================================

\subsection{Riemann Tensor}

For $\mathbb{C}P^n$ (complex dimension~$n$, real dimension $m = 2n$) with
holomorphic sectional curvature~$c$, the Riemann tensor in real indices is
\begin{equation}\label{eq:Riem-CPn}
R_{ijkl} = \frac{c}{4}\bigl[
  g_{ik}g_{jl} - g_{il}g_{jk}
  + J_{ik}J_{jl} - J_{il}J_{jk}
  + 2\,J_{ij}J_{kl}
\bigr],
\end{equation}
where $J$ is the complex structure (K\"ahler form as a $(1,1)$-tensor),
satisfying $J^2 = -\mathrm{Id}$, $J$ skew-symmetric, $\nabla J = 0$.

\subsection{Curvature Invariants}

For $\mathbb{C}P^2$ ($n = 2$, $m = 4$, $c = 4$):

\begin{proposition}[Ricci tensor and scalar curvature]
$\mathbb{C}P^2$ is Einstein:
\begin{equation}
\mathrm{Ric}_{ij} = \frac{(n+1)\,c}{2}\,g_{ij} = 6\,g_{ij},
\qquad
R = n(n+1)\,c = 24.
\end{equation}
\end{proposition}

\begin{proposition}[$|\mathrm{Ric}|^2$]
\begin{equation}
R_{ij}\,R^{ij} = \lambda^2 \cdot m = 36 \times 4 = 144.
\end{equation}
\end{proposition}

\begin{proposition}[Kretschner scalar $|{\rm Riem}|^2$]\label{prop:Kretschner}
Writing $R_{ijkl} = (c/4)\,T_{ijkl}$ with $T = A + B + C$ where
\begin{align}
A_{ijkl} &= g_{ik}g_{jl} - g_{il}g_{jk}, \\
B_{ijkl} &= J_{ik}J_{jl} - J_{il}J_{jk}, \\
C_{ijkl} &= 2\,J_{ij}J_{kl},
\end{align}
the inner products (computed in an orthonormal frame using
$J^TJ = \mathrm{Id}$, $J_{ii} = 0$) are:
\begin{center}
\begin{tabular}{lll}
\toprule
Inner product & General ($m = 2n$) & $m = 4$ \\
\midrule
$|A|^2$ & $2m(m-1)$ & 24 \\
$|B|^2$ & $2m(m-1)$ & 24 \\
$|C|^2$ & $4m^2$ & 64 \\
$\langle A, B\rangle$ & $2m$ & 8 \\
$\langle A, C\rangle$ & $4m$ & 16 \\
$\langle B, C\rangle$ & $4m$ & 16 \\
\bottomrule
\end{tabular}
\end{center}
Therefore:
\begin{equation}
|T|^2 = 24 + 24 + 64 + 16 + 32 + 32 = 192,
\end{equation}
\begin{equation}
\boxed{R_{ijkl}\,R^{ijkl} = \left(\frac{c}{4}\right)^{\!2} |T|^2 = 1 \times 192 = 192.}
\end{equation}
\end{proposition}

\subsection{Weyl Tensor Decomposition}

For a 4-dimensional Einstein manifold ($Z = \mathrm{Ric}_0 = 0$):
\begin{equation}
|\mathrm{Riem}|^2 = |W|^2 + \frac{R^2}{6},
\quad\Longrightarrow\quad
|W|^2 = 192 - \frac{576}{6} = 192 - 96 = 96.
\end{equation}
Since $\mathbb{C}P^2$ is self-dual ($W^- = 0$):
\begin{equation}
|W^+|^2 = 96, \qquad |W^-|^2 = 0.
\end{equation}

\subsection{Topological Invariants}

\begin{proposition}[Gauss--Bonnet verification]
Using $\mathrm{Vol}(\mathbb{C}P^2) = \pi^2/2$ and the 4D formula
\[
\chi(M^4) = \frac{1}{32\pi^2}\int_M
\bigl(|\mathrm{Riem}|^2 - 4\,|\mathrm{Ric}|^2 + R^2\bigr)\,\mathrm{d}\mathrm{vol},
\]
we obtain:
\[
\chi = \frac{1}{32\pi^2}\,(192 - 576 + 576)\,\frac{\pi^2}{2}
     = \frac{192}{64} = 3. \quad\checkmark
\]
\end{proposition}

\begin{proposition}[Hirzebruch signature verification]
Using
\[
\tau(M^4) = \frac{1}{48\pi^2}\int_M
\bigl(|W^+|^2 - |W^-|^2\bigr)\,\mathrm{d}\mathrm{vol},
\]
we obtain:
\[
\tau = \frac{96 \cdot \pi^2/2}{48\pi^2} = \frac{96}{96} = 1. \quad\checkmark
\]
\end{proposition}

\noindent
The first Pontryagin number is $p_1(\mathbb{C}P^2) = 3\tau = 3$.

\subsection{Summary: $\mathbb{C}P^2$}

\begin{center}
\begin{tabular}{ll}
\toprule
Quantity & Value \\
\midrule
$R$ & 24 \\
$\mathrm{Ric}_{ij} = \lambda\,g_{ij}$ & $\lambda = 6$ \\
$R_{ij}\,R^{ij}$ & 144 \\
$R_{ijkl}\,R^{ijkl}$ & 192 \\
$|W|^2$ & 96 \\
$|W^+|^2,\; |W^-|^2$ & 96, 0 \\
$\nabla^2 R$ & 0 \\
$\mathrm{Vol}$ & $\pi^2/2$ \\
$\chi$ & 3 \\
$\tau$ & 1 \\
$p_1$ & 3 \\
\bottomrule
\end{tabular}
\end{center}


%=============================================================================
\section{$S^3$ with Round Metric (Radius $r$)}
%=============================================================================

For $S^n$ of radius $r$ (constant sectional curvature $K = 1/r^2$):

\begin{align}
R_{ijkl} &= \frac{1}{r^2}\,(g_{ik}g_{jl} - g_{il}g_{jk}), \\
\mathrm{Ric}_{ij} &= \frac{n-1}{r^2}\,g_{ij} = \frac{2}{r^2}\,g_{ij}, \\
R &= \frac{n(n-1)}{r^2} = \frac{6}{r^2}.
\end{align}

\begin{center}
\begin{tabular}{ll}
\toprule
Quantity & Value ($n = 3$) \\
\midrule
$R$ & $6/r^2$ \\
$R_{ij}\,R^{ij}$ & $12/r^4$ \\
$R_{ijkl}\,R^{ijkl}$ & $12/r^4$ \\
$\nabla^2 R$ & 0 \\
$\mathrm{Vol}(S^3)$ & $2\pi^2 r^3$ \\
\bottomrule
\end{tabular}
\end{center}

\noindent
For $r = 1$: $R = 6$, $|{\rm Ric}|^2 = 12$, $|{\rm Riem}|^2 = 12$,
$\mathrm{Vol} = 2\pi^2$.


%=============================================================================
\section{Product Manifold $X = \mathbb{C}P^2 \times S^3$}
%=============================================================================

\subsection{Product Curvature Rules}

For a Riemannian product $M_1 \times M_2$ with block-diagonal metric
$g = g_1 \oplus g_2$:

\begin{lemma}[No cross terms in product Riemann tensor]
The Christoffel symbols of a product metric are block-diagonal. Therefore
the Riemann tensor vanishes whenever indices mix between $M_1$ and $M_2$:
\[
R^X_{ijkl} = \begin{cases}
  R^{M_1}_{ijkl} & \text{all indices in } M_1, \\
  R^{M_2}_{ijkl} & \text{all indices in } M_2, \\
  0 & \text{mixed indices.}
\end{cases}
\]
\end{lemma}

\noindent
\textbf{Consequences:}
\begin{align}
R_X &= R_1 + R_2, \\
|{\rm Ric}_X|^2 &= |{\rm Ric}_1|^2 + |{\rm Ric}_2|^2
  \quad\text{(block-diagonal Ricci)}, \\
|{\rm Riem}_X|^2 &= |{\rm Riem}_1|^2 + |{\rm Riem}_2|^2
  \quad\text{(\textbf{no} cross terms!)}.
\end{align}
However, $R_X^2 = (R_1 + R_2)^2$ \textbf{does} have a cross term
$2\,R_1\,R_2$.

\subsection{Numerical Values ($r = 1$)}

\begin{center}
\begin{tabular}{lccc}
\toprule
& $\mathbb{C}P^2$ & $S^3$ & $\mathbb{C}P^2 \times S^3$ \\
\midrule
$\dim$ & 4 & 3 & 7 \\
$R$ & 24 & 6 & 30 \\
$|{\rm Ric}|^2$ & 144 & 12 & 156 \\
$|{\rm Riem}|^2$ & 192 & 12 & 204 \\
$R^2$ & 576 & 36 & 900 \\
$\nabla^2 R$ & 0 & 0 & 0 \\
Vol (exact) & $\pi^2/2$ & $2\pi^2$ & $\pi^4$ \\
\bottomrule
\end{tabular}
\end{center}

\noindent
Note that $R_X^2 = 900 \neq 576 + 36 = 612$; the cross term is
$2 \times 24 \times 6 = 288$.


%=============================================================================
\section{The $a_4$ Seeley--DeWitt Coefficient}
%=============================================================================

\subsection{General Formula}

For a second-order operator $D = -(g^{ij}\nabla_i\nabla_j + E)$ acting on
sections of a vector bundle over a compact manifold without boundary, the
$a_4$ heat kernel coefficient density is~\cite{Vassilevich2003,Gilkey1995}:
\begin{equation}\label{eq:a4-general}
(4\pi)^{d/2}\,a_4(x) = \mathrm{tr}\Bigl[
  \tfrac{1}{2}E^2 + \tfrac{1}{6}\nabla^2 E
  + \tfrac{1}{12}\Omega_{ij}\Omega^{ij}
  + \tfrac{1}{30}\nabla^2 R
  + \tfrac{1}{72}R^2
  - \tfrac{1}{180}R_{ij}R^{ij}
  + \tfrac{1}{180}R_{ijkl}R^{ijkl}
  - \tfrac{1}{6}RE
\Bigr].
\end{equation}

\subsection{Minimally Coupled Real Scalar ($E = 0$, $\Omega = 0$)}

For a single minimally coupled real scalar field ($E = 0$, $\Omega = 0$,
$\mathrm{tr} = 1$), on a symmetric space ($\nabla^2 R = 0$):
\begin{equation}\label{eq:a4-minimal}
(4\pi)^{d/2}\,a_4(x) = \frac{1}{72}R^2
  - \frac{1}{180}|{\rm Ric}|^2
  + \frac{1}{180}|{\rm Riem}|^2.
\end{equation}

\begin{proposition}[$a_4$ density on $\mathbb{C}P^2 \times S^3$]
\begin{align}
(4\pi)^{7/2}\,a_4 &= \frac{1}{72}\times 900
  - \frac{1}{180}\times 156
  + \frac{1}{180}\times 204 \nonumber\\
&= \frac{25}{2} - \frac{13}{15} + \frac{17}{15} \nonumber\\
&= \frac{25}{2} + \frac{4}{15} \nonumber\\
&= \boxed{\frac{383}{30} \approx 12.7\overline{6}}\,.
\end{align}
\end{proposition}

\subsection{Integrated Coefficient}

The integrated $a_4$ coefficient is:
\begin{equation}
a_4[X] = \frac{\mathrm{Vol}(X)}{(4\pi)^{d/2}} \times a_4\text{-density}
= \frac{\pi^4}{(4\pi)^{7/2}} \times \frac{383}{30}.
\end{equation}
Numerically:
\begin{equation}
a_4[X] \approx 0.176784.
\end{equation}

\subsection{Dependence on $S^3$ Radius $r$}

With general $S^3$ radius $r$:
\begin{align}
R_X &= 24 + \frac{6}{r^2}, \\
|{\rm Ric}_X|^2 &= 144 + \frac{12}{r^4}, \\
|{\rm Riem}_X|^2 &= 192 + \frac{12}{r^4}.
\end{align}

The $a_4$ density becomes:
\begin{equation}\label{eq:a4-r}
(4\pi)^{7/2}\,a_4(r) = \frac{124}{15}
  + \frac{4}{r^2} + \frac{1}{2r^4}\,.
\end{equation}
The integrated coefficient is:
\begin{equation}
a_4[X](r) = \frac{\pi^4\,r^3}{(4\pi)^{7/2}}
  \left(\frac{124}{15} + \frac{4}{r^2} + \frac{1}{2r^4}\right).
\end{equation}

\subsection{Conformally Coupled Scalar ($d = 7$)}

For conformal coupling in dimension $d$:
\[
E = -\xi\,R, \qquad \xi = \frac{d-2}{4(d-1)} = \frac{5}{24}.
\]
The $R^2$ coefficient becomes:
\[
\frac{1}{2}\xi^2 + \frac{1}{6}\xi + \frac{1}{72}
= \frac{1}{2}\cdot\frac{25}{576} + \frac{1}{6}\cdot\frac{5}{24} + \frac{1}{72}
= \frac{9}{128} \approx 0.0703125.
\]
The $a_4$ density (conformal, $r = 1$) is approximately $63.548$,
giving an integrated $a_4 \approx 0.8800$.


%=============================================================================
\section{Individual Factor Checks}
%=============================================================================

\subsection{$a_4$ for $\mathbb{C}P^2$ Alone ($d = 4$)}

\begin{equation}
(4\pi)^2\,a_4 = \frac{576}{72} - \frac{144}{180} + \frac{192}{180}
= 8 - \frac{4}{5} + \frac{16}{15}
= \frac{124}{15} \approx 8.267.
\end{equation}
Integrated:
\begin{equation}
a_4[\mathbb{C}P^2] = \frac{\pi^2/2}{(4\pi)^2}\cdot\frac{124}{15}
= \frac{124}{480} = \frac{31}{120} \approx 0.2583.
\end{equation}

\subsection{Consistency Checks}

\begin{enumerate}
\item \textbf{Gauss--Bonnet:} $\chi(\mathbb{C}P^2) = 3$ verified from
$|{\rm Riem}|^2 - 4|{\rm Ric}|^2 + R^2 = 192$.
\item \textbf{Hirzebruch:} $\tau(\mathbb{C}P^2) = 1$ verified from
$|W^+|^2 - |W^-|^2 = 96$.
\item \textbf{Einstein decomposition:} $|{\rm Riem}|^2 = |W|^2 + R^2/6$
$\Rightarrow$ $192 = 96 + 96$. \checkmark
\item \textbf{Product Riemann:} No cross terms because Christoffel symbols
are block-diagonal for a product metric.
\end{enumerate}


%=============================================================================
\section{Key Results for the DFD $\alpha$ Derivation}
%=============================================================================

\begin{center}
\fbox{\parbox{0.9\textwidth}{%
\textbf{Summary of exact results} (all for $r = 1$):\smallskip

$X = \mathbb{C}P^2 \times S^3$, $\dim X = 7$.

\begin{tabular}{lrl}
Scalar curvature & $R_X ={}$ & $30$ \\
Ricci squared & $|{\rm Ric}_X|^2 ={}$ & $156$ \\
Kretschner scalar & $|{\rm Riem}_X|^2 ={}$ & $204$ \\
Volume & $\mathrm{Vol}(X) ={}$ & $\pi^4$ \\
$a_4$ density (minimal) & $(4\pi)^{7/2}\,a_4 ={}$ & $383/30$ \\
$a_4$ integrated (minimal) & $a_4[X] \approx{}$ & $0.1768$ \\
\end{tabular}

\medskip
The $a_4$ density as a function of $S^3$ radius:
\[
(4\pi)^{7/2}\,a_4(r) = \frac{124}{15} + \frac{4}{r^2} + \frac{1}{2r^4}\,.
\]
}}
\end{center}


\begin{thebibliography}{9}

\bibitem{Vassilevich2003}
D.\,V.~Vassilevich,
``Heat kernel expansion: user's manual,''
Phys.\ Rept.\ \textbf{388} (2003) 279--360,
\texttt{arXiv:hep-th/0306138}.

\bibitem{Gilkey1995}
P.\,B.~Gilkey,
\textit{Invariance Theory, the Heat Equation, and the Atiyah--Singer
Index Theorem},
2nd ed., CRC Press, 1995.

\bibitem{Besse1987}
A.\,L.~Besse,
\textit{Einstein Manifolds},
Springer-Verlag, 1987.

\end{thebibliography}

\end{document}
